\documentclass[portrait,final,a0paper,fontscale=0.277]{baposter}

\usepackage{calc}
\usepackage{graphicx}
\usepackage{amsmath}
\usepackage{amsthm}
\usepackage{amssymb}
\usepackage{relsize}
\usepackage{multirow}
\usepackage{rotating}
\usepackage{bm}
\usepackage{url}
%\usepackage{hyperref}  incompatiable with baposter
\usepackage{graphicx}
\usepackage{multicol}
\usepackage{pgfplots,pgfplotstable}
\usepackage{tikz}
\usetikzlibrary{plotmarks}
\pgfplotsset{compat=newest}
\usepackage{filecontents}
%\usepackage[perpage,para,symbol*]{footmisc}
%%%biblatex%%%%%%%%%%
%\usepackage[style=numeric,natbib=true]{biblatex}
%\usepackage[style=numeric]{biblatex}
%\addbibresource{/home/jim/JamesAshe/biblatex_jashe}
%\defbibheading{poster}{\vskip 0.1em} % Replace default heading
%\renewcommand{\bibfont}{\footnotesize\sffamily}


%\usepackage{times}
%\usepackage{helvet}
%\usepackage{bookman}
\usepackage{palatino}
%\usepackage{float}

\newcommand{\captionfont}{\footnotesize}
%\usepackage[font=small,labelfont=bf]{caption}


\graphicspath{{images/}{../images/}}
\usetikzlibrary{calc}

\newcommand{\SET}[1]  {\ensuremath{\mathcal{#1}}}
\newcommand{\MAT}[1]  {\ensuremath{\boldsymbol{#1}}}
\newcommand{\VEC}[1]  {\ensuremath{\boldsymbol{#1}}}
\newcommand{\Video}{\SET{V}}
\newcommand{\video}{\VEC{f}}
\newcommand{\track}{x}
\newcommand{\Track}{\SET T}
\newcommand{\LMs}{\SET L}
\newcommand{\lm}{l}
\newcommand{\PosE}{\SET P}
\newcommand{\posE}{\VEC p}
\newcommand{\negE}{\VEC n}
\newcommand{\NegE}{\SET N}
\newcommand{\Occluded}{\SET O}
\newcommand{\occluded}{o}
%%%theorem naming
\newtheoremstyle{named}{}{}{\itshape}{}{\bfseries}{.}{.5em}{\thmnote{#3 }#1} 
\theoremstyle{named} 
\newtheorem*{namedtheorem}{Theorem} 

%%%%%%%%%%%%%%%%%%%%%%%%%%%%%%%%%%%%%%%%%%%%%%%%%%%%%%%%%%%%%%%%%%%%%%%%%%%%%%%%
%%%% Some math symbols used in the text
%%%%%%%%%%%%%%%%%%%%%%%%%%%%%%%%%%%%%%%%%%%%%%%%%%%%%%%%%%%%%%%%%%%%%%%%%%%%%%%%

%%%%%%%%%%%%%%%%%%%%%%%%%%%%%%%%%%%%%%%%%%%%%%%%%%%%%%%%%%%%%%%%%%%%%%%%%%%%%%%%
% Multicol Settings
%%%%%%%%%%%%%%%%%%%%%%%%%%%%%%%%%%%%%%%%%%%%%%%%%%%%%%%%%%%%%%%%%%%%%%%%%%%%%%%%
\setlength{\columnsep}{1.5em}
\setlength{\columnseprule}{0mm}

%%%%%%%%%%%%%%%%%%%%%%%%%%%%%%%%%%%%%%%%%%%%%%%%%%%%%%%%%%%%%%%%%%%%%%%%%%%%%%%%
% Save space in lists. Use this after the opening of the list
%%%%%%%%%%%%%%%%%%%%%%%%%%%%%%%%%%%%%%%%%%%%%%%%%%%%%%%%%%%%%%%%%%%%%%%%%%%%%%%%
\newcommand{\compresslist}{%
\setlength{\itemsep}{1pt}%
\setlength{\parskip}{0pt}%
\setlength{\parsep}{0pt}%
}

%%%%%%%%%%%%%%%%%%%%%%%%%%%%%%%%%%%%%%%%%%%%%%%%%%%%%%%%%%%%%%%%%%%%%%%%%%%%%%
%%% Begin of Document
%%%%%%%%%%%%%%%%%%%%%%%%%%%%%%%%%%%%%%%%%%%%%%%%%%%%%%%%%%%%%%%%%%%%%%%%%%%%%%

\begin{document}

%%%%%%%%%%%%%%%%%%%%%%%%%%%%%%%%%%%%%%%%%%%%%%%%%%%%%%%%%%%%%%%%%%%%%%%%%%%%%%
%%% Here starts the poster
%%%---------------------------------------------------------------------------
%%% Format it to your taste with the options
%%%%%%%%%%%%%%%%%%%%%%%%%%%%%%%%%%%%%%%%%%%%%%%%%%%%%%%%%%%%%%%%%%%%%%%%%%%%%%
% Define some colors

%\definecolor{lightblue}{cmyk}{0.83,0.24,0,0.12}
\definecolor{lightblue}{rgb}{0.145,0.6666,1}
\definecolor{JCSUblue}{rgb}{0.0,0.32,1.59}
\definecolor{darkblue}{rgb}{0.0,0.32,2.0}
\definecolor{JCSUyellow}{rgb}{2.35,1.82,0.0}

\begin{poster}%
  % Poster Options
  {
  % Show grid to help with alignment
  grid=false,
  % Column spacing
  colspacing=1em,
  % Color style
  bgColorOne=white,
  bgColorTwo=white,
  borderColor=lightblue,
  headerColorOne=black,
  headerColorTwo=JCSUblue,
  headerFontColor=JCSUyellow,
  boxColorOne=white,
  boxColorTwo=JCSUblue,
  % Format of textbox
  textborder=roundedleft,
  % Format of text header
  eyecatcher=true,
  headerborder=closed,
  headerheight=0.1\textheight,
%  textfont=\sc, An example of changing the text font
  headershape=roundedright,
  headershade=shadelr,
  headerfont=\Large\bf\textsc, %Sans Serif
  textfont={\setlength{\parindent}{1.5em}},
  boxshade=plain,
%  background=shade-tb,
  background=plain,
  linewidth=2pt
  }
  % Eye Catcher
	{\includegraphics[height=5.5em]{images/Thomson_410_mesh_thermal.png}}
	%\captionfont{\emph{Energy distribution $N=410$ point}
  % Title
  {\bf\textsc{Modeling the Thomson Problem with Circle Packings}\vspace{0.5em}}
  % Authors
  {\textsc{ Daniel Villanueva  \hspace{1em} Merschia Griffin \hspace{1em} James Ashe, Ph.D. \\ \vspace{.1em} \normalsize{\textbf{Department of Natural Sciences and Mathematics, Johnson C. Smith University, Charlotte, NC}}}}
  % University logo
  {\includegraphics[height=5.5em]{images/jcsu_logo.png}
    %\includegraphics[width=5.5em]{images/jcsu_logo.jpg}}
  }
  

%%%%%%%%%%%%%%%%%%%%%%%%%%%%%%%%%%%%%%%%%%%%%%%%%%%%%%%%%%%%%%%%%%%%%%%%%%%%%%
%%% Now define the boxes that make up the poster
%%%---------------------------------------------------------------------------
%%% Each box has a name and can be placed absolutely or relatively.
%%% The only inconvenience is that you can only specify a relative position 
%%% towards an already declared box. So if you have a box attached to the 
%%% bottom, one to the top and a third one which should be in between, you 
%%% have to specify the top and bottom boxes before you specify the middle 
%%% box.
%%%%%%%%%%%%%%%%%%%%%%%%%%%%%%%%%%%%%%%%%%%%%%%%%%%%%%%%%%%%%%%%%%%%%%%%%%%%%%
    %
    % A coloured circle useful as a bullet with an adjustably strong filling
    \newcommand{\colouredcircle}{%
      \tikz{\useasboundingbox (-0.2em,-0.32em) rectangle(0.2em,0.32em); \draw[draw=black,fill=lightblue,line width=0.03em] (0,0) circle(0.18em);}}

%%%%%%%%%%%%%%%%%%%%%%%%%%%%%%%%%%%%%%%%%%%%%%%%%%%%%%%%%%%%%%%%%%%%%%%%%%%%%%
  \headerbox{The Thomson Problem}{name=problem,column=0,row=0,span=1}{
%%%%%%%%%%%%%%%%%%%%%%%%%%%%%%%%%%%%%%%%%%%%%%%%%%%%%%%%%%%%%%%%%%%%%%%%%%%%%%
The Thomson problem asks how to find the minimum energy distribution of $N$ electrons on the surface of the sphere. The \emph{energy}
of $N$ points, $p_{1},p_{2},\dots,p_{N}$, on the sphere is defined
as 
\[
E={\displaystyle \sum_{i<j}\frac{1}{\left|p_{i}-p_{j}\right|}}\mbox{.}
\]
Number 7 on Smale's list of conjectures, this problem remains one
of the most important unsolved packing problems, and has applications in biology and physics \cite{wiki}.
%Not just important
%in mathematics, the Thomson problem is related or equivalent to many
%other problems in biology, physics, computer science, chemisty, etc
%{[}wiki{]}. 

Known solutions have been rigorously proven for only a small number
of special cases. In general, finding a solution is an NP hard problem, and progress has relied on computational methods
\cite{wales}.
\\[-.8em]
%Using such methods, smallest known energy distributions for most points
%sets less than 1000 have been found {[}cite paper{]}.
}

%%%%%%%%%%%%%%%%%%%%%%%%%%%%%%%%%%%%%%%%%%%%%%%%%%%%%%%%%%%%%%%%%%%%%%%%%%%%%%
  \headerbox{Circle Packings}{name=packing,column=0,below=problem, above=bottom}{
%%%%%%%%%%%%%%%%%%%%%%%%%%%%%%%%%%%%%%%%%%%%%%%%%%%%%%%%%%%%%%%%%%%%%%%%%%%%%%
   Circle packings are configurations of circles with specified patterns
of tangency. The following theorem provides rigidity and existence
of circle packings on the sphere.

\begin{namedtheorem}[Koebe-Andreev-Thurston]

There exists a circle packing on the sphere for any finite triangulation
of the sphere. Furthermore, this circle packing is unique up to M\"{o}bius
transformations.
\end{namedtheorem} 

\vspace{-.5em} 
\noindent\begin{tabular}{@{\hspace{0.0em}}c@{\hspace{0.0em}}c@{\hspace{0.0em}}}
\includegraphics[width=1.0\linewidth]{images/complex_to_sphere.pdf}\\[-.5em]
\captionfont{\emph{A tangency pattern and its circle packing.}}
\end{tabular}
\\[.5em]
The actual radii for the circle packing guaranteed by this theorem
can always be found using the \emph{Circle Packing Algorithm} \cite{algor}.
\renewcommand{\thefootnote}{\fnsymbol{footnote}}
Moreover we have a feature rich software package,  $\tt{CirclePack}$\footnote{\url{http://www.math.utk.edu/~kens/CirclePack/};
special thanks to Ken Stephenson, University of Tennessee, for developing and maintaining this
freely available software.
}, for computing and visualizing the lovely results.

\begin{namedtheorem}[Thurston's Circle Packing Conjecture]

The conformal mapping of a bounded simply connected domain onto the
unit disc can be approximated with a circle packing.

\end{namedtheorem}

In summary, we have the following:
\begin{enumerate}
\item An essentially unique circle packing for a spherical triangulation
exists.
\item The circle packing can be computationally realized using $\tt{CirclePack}$.
\item As the number of circles increase, a circle packing analogue will
converge towards its continuous counter part  \cite{introCP}. 
\end{enumerate}
}

%%%%%%%%%%%%%%%%%%%%%%%%%%%%%%%%%%%%%%%%%%%%%%%%%%%%%%%%%%%%%%%%%%%%%%%%%%%%%%
\headerbox{Combinatorics}{name=combinatorics,column=2,row=0}{
%%%%%%%%%%%%%%%%%%%%%%%%%%%%%%%%%%%%%%%%%%%%%%%%%%%%%%%%%%%%%%%%%%%%%%%%%%%%%%
   Repelling points on a flat surface naturally adopt a hexagonal pattern. However the topology of the sphere precludes a triangulation where every vertex has 6 neighbors. Try wrapping a sheet of paper around a basketball --imperfections in your wrapping job are unavoidable.

The Euler characteristic states that there will be 12 such ``imperfections'' where a vertex has 5 neighbors. The truncated icosahedron, a soccer ball, is a classic example. Thus a combinatorial perspective may shed new light on this problem.  
}

%%%%%%%%%%%%%%%%%%%%%%%%%%%%%%%%%%%%%%%%%%%%%%%%%%%%%%%%%%%%%%%%%%%%%%%%%%%%%%
\headerbox{Manipulations}{name=manipulations,column=2,below=combinatorics,row=0}{
%%%%%%%%%%%%%%%%%%%%%%%%%%%%%%%%%%%%%%%%%%%%%%%%%%%%%%%%%%%%%%%%%%%%%%%%%%%%%%
Now it is demonstrated how manipulating the combinatorics of circle packings can be used to reduce energy.      
    A Whitehead move, or ``flip'', replaces one interior edge with another in the union of two edge-sharing faces. Such a modification does not change the number of points -only the combinatorics locally. However these local changes result in a new circle packing and thus invoke a global geometric change.
    
    A common issue with optimization algorithms is settling into a local but not a global minimum. In the figures below, two defects are separated be an area of degree six vertices. Following a sequence of flips, the defects are then combined and neutralized. Locally such defects may be difficult to fix, but the ``ladder'' bridges the two local imperfections.
    \\[.25em]
    \noindent\begin{tabular}{@{\hspace{0.0em}}c@{\hspace{0.0em}}c@{\hspace{0.0em}}}
	\includegraphics[width=1.0\linewidth]{images/ladder.pdf}\\[-.5em]
	\captionfont{\emph{A defect is eliminated through a series of flips.}}\\
	\end{tabular}	
}

%%%%%%%%%%%%%%%%%%%%%%%%%%%%%%%%%%%%%%%%%%%%%%%%%%%%%%%%%%%%%%%%%%%%%%%%%%%%%%
\headerbox{A Discrete Model}{name=model,column=1,row=0}{
%%%%%%%%%%%%%%%%%%%%%%%%%%%%%%%%%%%%%%%%%%%%%%%%%%%%%%%%%%%%%%%%%%%%%%%%%%%%%%
The Thomson problem is well studied, and there is a large catalog
of current minimal energy distributions \cite{thomson:wales}. These point sets are stored
in the interactive web applet, $\tt{thomsonapplet}$\footnote{\url{http://thomson.phy.syr.edu}, Cris
Cecka, Primary Author. ccecka@seas.harvard.edu}, which can also generate random point sets, measure the energy, and
implement a wide array of algorithms. 

To construct a circle packing model for $N$ energy points on the
sphere we did the following: 
\begin{enumerate}
\item Find the Delaunay triangulations for $N$ points using data from $\tt{thomsonapplet}$. 
\item Compute the circle packing for this triangulation using $\tt{CirclePack}$.
\item Normalize the packing by applying a M\"{o}bius map such that the
centroid of the circle centers in $\mathbb{R}^{3}$ is at the origin. 
\end{enumerate}
The \emph{error }of this model was measured as, 
\[
error={\captionfont \frac{\mbox{(energy from circle packing)-(energy)}}{\mbox{energy}}}
\]
\\[-1.25em]
\noindent\begin{tabular}{@{\hspace{0.0em}}c@{\hspace{0.0em}}c@{\hspace{0.0em}}}
\includegraphics[width=1.0\linewidth]{images/data_and_chart_for_printing3.pdf}
\\[-.5em]
\end{tabular}
The Figure above shows the error for minimal known $N$ distributions. The error from our experiments never exceeded $0.0062$, and as
expected the error tended toward $0$ for large values of $N$. The error introduced by the discreteness of circle packings
is limited for even very small $N$ and inconsequential for large $N$.

\noindent\begin{tabular}{@{\hspace{0.0em}}c@{\hspace{1.5em}}c@{\hspace{0.0em}}}
%\includegraphics[width=0.46\linewidth]{images/Thomson_410_mesh_thermal.png} &
\includegraphics[width=0.46\linewidth]{images/Thomson_410_3d_pts.png} &
\includegraphics[width=0.46\linewidth]{images/cp410_min.png} \\
\captionfont{\emph{A minimal distributon}} & \captionfont{\emph{A circle packing model}} 
\end{tabular}
}

%%%%%%%%%%%%%%%%%%%%%%%%%%%%%%%%%%%%%%%%%%%%%%%%%%%%%%%%%%%%%%%%%%%%%%%%%%%%%%
\headerbox{Acknowledgments}{name=thanks,column=2,row=0,below=manipulations,above=bottom}{
%%%%%%%%%%%%%%%%%%%%%%%%%%%%%%%%%%%%%%%%%%%%%%%%%%%%%%%%%%%%%%%%%%%%%%%%%%%%%%
\smaller
Funding for this project was provided by HBCU-UP of Johnson C.\,Smith. Ken Stephenson, University of Tennessee, provided advice and programming support. $\tt{CirclePack}$, $\tt{thomsonapplet}$, $\tt{Inkscape}$, $\tt{Gimp}$, and \LaTeX{} were all used to produce this document.
}

%%%%%%%%%%%%%%%%%%%%%%%%%%%%%%%%%%%%%%%%%%%%%%%%%%%%%%%%%%%%%%%%%%%%%%%%%%%%%%
  \headerbox{References}{name=references,column=1,below=model,above=bottom}{
%%%%%%%%%%%%%%%%%%%%%%%%%%%%%%%%%%%%%%%%%%%%%%%%%%%%%%%%%%%%%%%%%%%%%%%%%%%%%%
    \smaller
    \bibliographystyle{ieee}
    \renewcommand{\section}[2]{\vskip 0.05em}
    \begin{thebibliography}{1}\itemsep=-0.01em
      \setlength{\baselineskip}{0.4em}
      
      \bibitem{wiki}
	\newblock \emph{Thomson Problem} 2012.
	\newblock \url{http://en.wikipedia.org/wiki/Thomson_problem}.
      
      \bibitem{wales}
	D.~J.~Wales and S.~Ulker.
	\newblock ``Structure and dynamics of spherical crystals characterized for the Thomson problem''.
	\newblock \emph{Physical Review B} 74.21 (2006).

      \bibitem{algor}
        C.~Collins and K.~Stephenson.
        \newblock ``A circle packing algorithm''.
	\newblock \emph{Computational Geometry} 25.3 (2003).

      \bibitem{introCP}
	K.~Stephenson.
	\newblock \emph{Introduction to circle packing}. 
	\newblock Cambridge University Press (2005). 

      \bibitem{thomson:wales}
	D.~J.~Wales and S.~Ulker.
	\newblock \emph{Global Minima for the Thomson Problem} (2006). 
	\newblock \url{ http://www-							wales.ch.cam.ac.uk/~wales/CCD/Thomson/table.html}.
    \end{thebibliography}
  \vspace{0.3em}
  }
  
\end{poster}

\end{document}
